\chapter{Symmetries, Noether Currents, and Stress Tensors}
\label{ch:symmetries-noether-stress}
\ConventionsLorentzian


\section{Variational Framework}

\begin{definition}[First-derivative classical field datum]
\label{def:first-derivative-classical-field-datum}
Work on \(D\)-dimensional Minkowski spacetime \(\mathbb M^D\), with
mostly-plus metric \(\eta_{\mu\nu}\).  Let \(A\) range over a finite set of
bosonic field components.  A classical field configuration is a smooth map
\[
  \phi:M\longrightarrow \mathbb R^N,\qquad
  x\longmapsto \{\phi^A(x)\},
\]
on a spacetime region \(M\subset \mathbb M^D\).  In this chapter a Lagrangian
density means a first-derivative local expression
\[
  \mathcal L=\mathcal L(\phi^A,\partial_\mu\phi^A)
\]
with no dependence on derivatives of order two or higher.  The corresponding
action over \(M\) is
\[
  S_M[\phi]=\int_M\dd^Dx\,\mathcal L(\phi,\partial\phi).
\]
A variation is a smooth one-parameter family \(\phi_s\) with
\(\phi_0=\phi\); its tangent is
\[
  \delta\phi^A=\left.\frac{\dd}{\dd s}\phi_s^A\right|_{s=0}.
\]
When the boundary term is not part of the question, \(\delta\phi^A\) is
assumed to have compact support in \(M^\circ\), the interior of \(M\).
\end{definition}

\begin{proposition}[First variational identity]
\label{prop:first-variational-identity}
For the datum of
Definition~\ref{def:first-derivative-classical-field-datum}, define the
Euler--Lagrange expression
\[
  E_A(\mathcal L)
  =
  \frac{\partial\mathcal L}{\partial\phi^A}
  -
  \partial_\mu
  \frac{\partial\mathcal L}{\partial(\partial_\mu\phi^A)} .
\]
Then every smooth variation satisfies
\[
  \delta\mathcal L
  =
  E_A(\mathcal L)\,\delta\phi^A
  +
  \partial_\mu\theta^\mu(\delta\phi),
  \qquad
  \theta^\mu(\delta\phi)
  =
  \frac{\partial\mathcal L}{\partial(\partial_\mu\phi^A)}
  \delta\phi^A ,
\]
where repeated component and spacetime indices are summed.  Consequently
\[
  \delta S_M
  =
  \int_M\dd^Dx\,E_A(\mathcal L)\,\delta\phi^A
  +
  \int_{\partial M}\dd\Sigma_\mu\,\theta^\mu(\delta\phi),
\]
where \(\dd\Sigma_\mu\) is the oriented hypersurface measure on
\(\partial M\).
\end{proposition}

\begin{proof}
The chain rule gives
\[
  \delta\mathcal L
  =
  \frac{\partial\mathcal L}{\partial\phi^A}\delta\phi^A
  +
  \frac{\partial\mathcal L}{\partial(\partial_\mu\phi^A)}
  \partial_\mu\delta\phi^A .
\]
Apply the product rule to the second term:
\[
  \frac{\partial\mathcal L}{\partial(\partial_\mu\phi^A)}
  \partial_\mu\delta\phi^A
  =
  \partial_\mu
  \left(
    \frac{\partial\mathcal L}{\partial(\partial_\mu\phi^A)}
    \delta\phi^A
  \right)
  -
  \partial_\mu
  \frac{\partial\mathcal L}{\partial(\partial_\mu\phi^A)}
  \delta\phi^A .
\]
Combining the two non-divergence terms gives \(E_A(\mathcal L)\delta\phi^A\).
Integrating the pointwise identity over \(M\) and applying Stokes' theorem
gives the displayed formula for \(\delta S_M\).
\end{proof}

\begin{definition}[Classical solution]
\label{def:classical-solution-euler-lagrange}
A classical solution of the first-derivative Lagrangian problem is a field
configuration obeying
\[
  E_A(\mathcal L)=0
\]
for every component \(A\), together with the boundary or asymptotic conditions
specified by the variational problem.  The Euler--Lagrange equation alone is
not a complete boundary-value or Cauchy problem; the boundary data are part of
the definition of the classical problem being solved.
\end{definition}

\section{Infinitesimal Symmetries}

\begin{definition}[Infinitesimal variational symmetry]
\label{def:infinitesimal-variational-symmetry}
An infinitesimal field-space transformation is a rule
\[
  \delta_R\phi^A=R^A[\phi],
\]
where \(R^A[\phi]\) is a local expression in \(x\), the fields, and finitely
many derivatives of the fields.  The transformation is a variational symmetry
of the action if there exists a local vector \(K_R^\mu\) such that
\[
  \delta_R\mathcal L=\partial_\mu K_R^\mu .
\]
This equation is the definition of the symmetry in the present classical
Lagrangian framework.  It is an off-shell identity of local functions of the
fields; no Euler--Lagrange equation has yet been imposed.
\end{definition}

\begin{proposition}[Noether identity for a variational symmetry]
\label{prop:noether-identity-variational-symmetry}
Let \(R\) be a variational symmetry in the sense of
Definition~\ref{def:infinitesimal-variational-symmetry}.  The local current
\[
  j_R^\mu
  =
  \frac{\partial\mathcal L}{\partial(\partial_\mu\phi^A)}R^A
  -
  K_R^\mu
\]
obeys the off-shell identity
\[
  E_A(\mathcal L)R^A+\partial_\mu j_R^\mu=0.
\]
In particular, on every classical solution,
\[
  \partial_\mu j_R^\mu=0 .
\]
\end{proposition}

\begin{proof}
Apply Proposition~\ref{prop:first-variational-identity} to the variation
\(\delta_R\phi^A=R^A\):
\[
  \delta_R\mathcal L
  =
  E_A(\mathcal L)R^A
  +
  \partial_\mu
  \left(
    \frac{\partial\mathcal L}{\partial(\partial_\mu\phi^A)}R^A
  \right).
\]
The variational-symmetry condition says that the left hand side is
\(\partial_\mu K_R^\mu\).  Moving this total derivative to the right gives the
displayed off-shell identity.  Setting \(E_A(\mathcal L)=0\) gives current
conservation on solutions.
\end{proof}

\begin{proposition}[Localized-parameter form of Noether's theorem]
\label{prop:localized-parameter-noether}
Let \(\epsilon\in C_c^\infty(M)\), and localize the transformation by
\[
  \delta_{\epsilon R}\phi^A=\epsilon R^A .
\]
For a first-derivative Lagrangian,
\[
  \delta_{\epsilon R}\mathcal L
  =
  \partial_\mu(\epsilon K_R^\mu)
  +
  (\partial_\mu\epsilon)j_R^\mu .
\]
Thus, for compactly supported \(\epsilon\),
\[
  \delta_{\epsilon R}S_M
  =
  \int_M\dd^Dx\,(\partial_\mu\epsilon)\,j_R^\mu .
\]
If the field is a classical solution, this identity is equivalent, as a
distributional statement, to \(\partial_\mu j_R^\mu=0\).
\end{proposition}

\begin{proof}
Using the chain rule,
\[
  \delta_{\epsilon R}\mathcal L
  =
  \epsilon\,\delta_R\mathcal L
  +
  (\partial_\mu\epsilon)
  \frac{\partial\mathcal L}{\partial(\partial_\mu\phi^A)}R^A .
\]
Since \(\delta_R\mathcal L=\partial_\mu K_R^\mu\),
\[
  \epsilon\,\partial_\mu K_R^\mu
  =
  \partial_\mu(\epsilon K_R^\mu)-(\partial_\mu\epsilon)K_R^\mu .
\]
Combining the two \(\partial_\mu\epsilon\) terms gives
\((\partial_\mu\epsilon)j_R^\mu\).  The integral of
\(\partial_\mu(\epsilon K_R^\mu)\) vanishes because \(\epsilon\) has compact
support in \(M\).  On a solution the first variation of the action by a
compactly supported variation vanishes, hence
\[
  0
  =
  \int_M\dd^Dx\,(\partial_\mu\epsilon)j_R^\mu
  =
  -\int_M\dd^Dx\,\epsilon\,\partial_\mu j_R^\mu .
\]
The fundamental lemma of distributions then gives
\(\partial_\mu j_R^\mu=0\).  The converse follows by integrating by parts.
\end{proof}

\begin{definition}[Current improvement]
\label{def:current-improvement}
The Noether current is a representative of a local equivalence class.  If
\[
  B_R^{\nu\mu}=-B_R^{\mu\nu},
\]
then
\[
  j_R^\mu\longmapsto j_R^\mu+\partial_\nu B_R^{\nu\mu}
\]
is called a current improvement.  It has the same local divergence because
\(\partial_\mu\partial_\nu B_R^{\nu\mu}=0\).  The corresponding integrated
charge is unchanged when the surface integral of \(B_R\) at spatial boundary
vanishes in the stated boundary condition or asymptotic domain.
\end{definition}

\section{Charges and Generators}

\begin{definition}[Classical charge of a conserved current]
\label{def:classical-current-charge}
Let \(\Sigma_t=\{x^0=t\}\) be a constant-time slice.  If \(j_R^\mu\) is
conserved and the spatial integral exists, the charge on \(\Sigma_t\) is
\[
  Q_R(t)=\int_{\Sigma_t}\dd^{D-1}\vec x\,j_R^0(t,\vec x).
\]
For a spatial region \(V\subset\mathbb R^{D-1}\), the finite-region charge is
\[
  Q_R^V(t)=\int_V\dd^{D-1}\vec x\,j_R^0(t,\vec x).
\]
The existence of \(Q_R(t)\) is an analytic condition on the solution and on
the asymptotic domain; it is not implied by the local equation
\(\partial_\mu j_R^\mu=0\) alone.
\end{definition}

\begin{proposition}[Charge balance in a spacetime slab]
\label{prop:charge-balance-spacetime-slab}
Let \(V\subset\mathbb R^{D-1}\) be a spatial region with piecewise smooth
boundary, and assume \(j_R^\mu\) is smooth on
\([t_1,t_2]\times V\) with \(\partial_\mu j_R^\mu=0\).  Then
\[
  Q_R^V(t_2)-Q_R^V(t_1)
  =
  -\int_{t_1}^{t_2}\dd x^0
  \int_{\partial V}\dd S_i\,j_R^i(x^0,\vec x),
\]
where \(\dd S_i=n_i\,\dd S\), \(n_i\) is the outward unit normal to
\(\partial V\), and \(\dd S\) is the Euclidean surface measure on
\(\partial V\).  Hence \(Q_R(t)\) is time independent whenever the boundary
flux tends to zero in the limit \(V\nearrow\mathbb R^{D-1}\).
\end{proposition}

\begin{proof}
Integrate \(\partial_0 j_R^0+\partial_i j_R^i=0\) over
\([t_1,t_2]\times V\).  The time derivative integrates to
\(Q_R^V(t_2)-Q_R^V(t_1)\), and the spatial divergence integrates by the
ordinary divergence theorem to
\(\int_{t_1}^{t_2}\dd x^0\int_{\partial V}\dd S_i\,j_R^i\).  Moving this
boundary term to the other side gives the stated identity.
\end{proof}

Figure~\ref{fig:current-conservation-boundary-flux} records the slab geometry
behind the charge-conservation identity: time-slice charges differ exactly by
flux through the spatial boundary.

\begin{figure}[H]
\centering
\begin{tikzpicture}[scale=0.95,>=Stealth]
  \draw[->,line width=0.7pt] (0,-0.25) -- (0,4.1) node[above] {$x^0$};
  \draw[->,line width=0.7pt] (-0.35,0) -- (5.4,0) node[right] {space};
  \draw[blue!65!black,line width=1pt] (1.05,0.8) -- (4.25,0.8);
  \draw[blue!65!black,line width=1pt] (1.05,3.15) -- (4.25,3.15);
  \draw[gray!70,line width=0.9pt] (1.05,0.8) -- (1.05,3.15);
  \draw[gray!70,line width=0.9pt] (4.25,0.8) -- (4.25,3.15);
  \node[blue!65!black,anchor=east] at (0.95,0.8) {$\Sigma_{t_1}$};
  \node[blue!65!black,anchor=east] at (0.95,3.15) {$\Sigma_{t_2}$};
  \node at (2.65,2.05) {$\partial_\mu j^\mu_R=0$};
  \draw[orange!80!black,thick,->] (4.25,2.05) -- (5.0,2.05)
    node[right] {boundary flux};
  \draw[green!45!black,thick,->] (2.65,0.8) -- (2.65,1.3);
  \draw[green!45!black,thick,->] (2.65,3.15) -- (2.65,2.65);
\end{tikzpicture}
\caption{Current conservation in a spacetime slab.  The difference between
charges on \(\Sigma_{t_1}\) and \(\Sigma_{t_2}\) is the flux through the
spatial boundary.}
\label{fig:current-conservation-boundary-flux}
\end{figure}

\begin{definition}[Equal-time Poisson bracket and Hamiltonian generator]
\label{def:equal-time-poisson-generator}
Let \(\Pi_A\) denote the canonical momentum conjugate to \(\phi^A\).  On the
phase-space functionals whose variational derivatives and boundary terms are
defined, the equal-time Poisson bracket is
\[
  \{F,G\}
  =
  \int\dd^{D-1}\vec x
  \left[
    \frac{\delta F}{\delta\phi^A(\vec x)}
    \frac{\delta G}{\delta\Pi_A(\vec x)}
    -
    \frac{\delta F}{\delta\Pi_A(\vec x)}
    \frac{\delta G}{\delta\phi^A(\vec x)}
  \right].
\]
In particular,
\[
  \{\phi^A(\vec x),\Pi_B(\vec y)\}
  =
  \delta^A{}_B\,\delta^{(D-1)}(\vec x-\vec y).
\]
A charge \(Q_R\) generates the infinitesimal transformation \(R\) on canonical
observables if
\[
  \delta_R F=\{F,Q_R\}
\]
for every functional \(F\) in the stated domain.  Applied to the coordinate
field this condition reads
\[
  R^A(t,\vec x)=\{\phi^A(t,\vec x),Q_R\}.
\]
\end{definition}

\begin{proposition}[First-order Noether charges as phase-space generators]
\label{prop:first-order-noether-charge-generator}
Assume \(K_R^0=0\), and assume \(R^A\) depends on the canonical coordinates
and their spatial derivatives but not on the velocities or momenta.  If the
Noether charge is represented on phase space by
\[
  Q_R=\int_{\Sigma_t}\dd^{D-1}\vec x\,\Pi_A R^A
\]
up to spatial-divergence improvements whose boundary integrals vanish, then
\[
  \{\phi^A(t,\vec x),Q_R\}=R^A(t,\vec x).
\]
More generally, if \(R^A\) or \(K_R^0\) contains velocities, the current
obtained from the Lagrangian calculation must first be expressed as a
functional on phase space before the generator statement is tested.
\end{proposition}

\begin{proof}
Under the stated hypotheses the functional derivative of \(Q_R\) with respect
to \(\Pi_A(\vec x)\) is \(R^A(\vec x)\).  The equal-time bracket in
Definition~\ref{def:equal-time-poisson-generator} therefore gives
\[
  \{\phi^A(\vec x),Q_R\}
  =
  \frac{\delta Q_R}{\delta\Pi_A(\vec x)}
  =
  R^A(\vec x).
\]
Spatial-divergence improvements do not change the result under the assumed
boundary condition.  If velocities occur, the expression
\(\int j_R^0\) is not yet a function on canonical phase space; the Legendre
map must be used first, and the displayed derivative calculation is then
performed on the resulting phase-space functional.
\end{proof}

\begin{example}[Time translation]
\label{ex:time-translation-noether-generator}
For a time translation by \(\tau\),
\[
  R^A=-\tau\dot\phi^A,\qquad Q_R=-\tau H.
\]
Hamilton's equation yields
\[
  \{\phi^A,-\tau H\}=-\tau\dot\phi^A .
\]
Thus the generator statement is a phase-space statement with the Noether
charge expressed in canonical variables.
\end{example}

\section{Poincare Transformations of Scalar Fields}

\begin{definition}[Scalar representation of the Poincare group]
\label{def:scalar-poincare-transformation}
For a scalar field, a finite proper orthochronous Poincare transformation is
described by
\[
  x'^\mu=\Lambda^\mu{}_\nu x^\nu+a^\mu,\qquad
  \phi'(x')=\phi(x),
\]
where \(\Lambda^T\eta\Lambda=\eta\).  Equivalently,
\[
  \phi'(x)
  =
  \phi\big(\Lambda^{-1}(x-a)\big).
\]
Write
\(\Lambda^\mu{}_\nu=\delta^\mu{}_\nu+\omega^\mu{}_\nu+O(\omega^2)\), with
\(\omega_{\mu\nu}=-\omega_{\nu\mu}\).  The fixed-coordinate infinitesimal
variation is
\[
  \delta\phi(x)
  =
  -\left(a^\mu+\omega^\mu{}_\nu x^\nu\right)\partial_\mu\phi(x).
\]
The translation and Lorentz currents below are the Noether currents for the
two independent parameters \(a^\mu\) and \(\omega_{\mu\nu}\).
\end{definition}

\section{Translation Symmetry and the Stress Tensor}

\begin{proposition}[Canonical stress tensor from translations]
\label{prop:canonical-stress-tensor-translations}
Assume that \(\mathcal L\) has no explicit dependence on \(x\).  For a
constant vector \(\xi^\nu\), use the fixed-coordinate field variation
\[
  \delta_\xi\phi^A(x)=-\xi^\nu\partial_\nu\phi^A(x).
\]
Then the translation Noether current can be written
\[
  j^\mu_\xi=\xi^\nu T^\mu{}_\nu,
\]
where
\[
  T^\mu{}_\nu
  =
  -
  \frac{\partial\mathcal L}{\partial(\partial_\mu\phi^A)}
  \partial_\nu\phi^A
  +
  \delta^\mu{}_\nu\,\mathcal L .
\]
On solutions,
\[
  \partial_\mu T^\mu{}_\nu=0.
\]
Equivalently, for compactly supported \(\xi^\nu(x)\),
\[
  \delta_\xi S_M
  =
  \int_M\dd^Dx\,(\partial_\mu\xi^\nu)\,T^\mu{}_\nu ,
\]
after the compactly supported total derivative has been separated.  Thus
\(T^\mu{}_\nu\) is the response coefficient for a localized translation.
\end{proposition}

\begin{proof}
For constant \(\xi^\nu\), the density has no explicit coordinate dependence.
Translation invariance of the scalar local density gives
\[
  \delta_\xi\mathcal L=-\xi^\nu\partial_\nu\mathcal L
  =
  \partial_\mu K_\xi^\mu,
  \qquad
  K_\xi^\mu=-\xi^\mu\mathcal L .
\]
Substituting \(R^A=-\xi^\nu\partial_\nu\phi^A\) and this \(K_\xi^\mu\) into
Proposition~\ref{prop:noether-identity-variational-symmetry} gives
\[
  j_\xi^\mu
  =
  -\xi^\nu
  \frac{\partial\mathcal L}{\partial(\partial_\mu\phi^A)}
  \partial_\nu\phi^A
  +
  \xi^\mu\mathcal L
  =
  \xi^\nu T^\mu{}_\nu .
\]
Since \(\xi^\nu\) is arbitrary, conservation of \(j_\xi^\mu\) for every
constant \(\xi^\nu\) is equivalent to
\(\partial_\mu T^\mu{}_\nu=0\).

For local \(\xi^\nu(x)\), compute directly:
\[
  \delta_\xi\mathcal L
  =
  -\xi^\nu\partial_\nu\mathcal L
  -
  (\partial_\mu\xi^\nu)
  \frac{\partial\mathcal L}{\partial(\partial_\mu\phi^A)}
  \partial_\nu\phi^A .
\]
The first term is
\(-\partial_\mu(\xi^\mu\mathcal L)+(\partial_\mu\xi^\mu)\mathcal L\).
The coefficient of \(\partial_\mu\xi^\nu\) is therefore precisely
\(T^\mu{}_\nu\).  Integration over \(M\) gives the localized-translation
formula because the total derivative has compact support.
\end{proof}

\begin{definition}[Translation charges]
\label{def:translation-charges}
When the spatial integrals exist, the momentum covector and vector are
\[
  P_\nu=\int\dd^{D-1}\vec x\,T^0{}_\nu,
  \qquad
  P^\nu=\eta^{\nu\rho}P_\rho
  =
  \int\dd^{D-1}\vec x\,T^{0\nu}.
\]
With the mostly-plus metric, \(P^0\) is the energy.  Equivalently,
\(P_0=-P^0\).
\end{definition}

\begin{proposition}[Stress tensor of a real scalar field]
\label{prop:real-scalar-stress-tensor}
For
\[
  \mathcal L
  =
  -\frac12\partial_\rho\phi\,\partial^\rho\phi
  -
  V(\phi),
\]
the canonical stress tensor is
\[
  T^\mu{}_\nu
  =
  \partial^\mu\phi\,\partial_\nu\phi
  -
  \delta^\mu{}_\nu
  \left(
    \frac12\partial_\rho\phi\,\partial^\rho\phi+V(\phi)
  \right).
\]
Raising the second index,
\[
  T^{\mu\nu}=\eta^{\nu\rho}T^\mu{}_\rho,
\]
gives a symmetric tensor in this example.  Its energy density is
\[
  T^{00}
  =
  \frac12(\partial_0\phi)^2
  +
  \frac12\sum_{i=1}^{D-1}(\partial_i\phi)^2
  +
  V(\phi).
\]
For the free scalar field, \(V(\phi)=\frac12m^2\phi^2\), this is the
Hamiltonian density obtained by canonical quantization.
\end{proposition}

\begin{proof}
The derivative of the Lagrangian with respect to \(\partial_\mu\phi\) is
\[
  \frac{\partial\mathcal L}{\partial(\partial_\mu\phi)}
  =
  -\partial^\mu\phi .
\]
Substitution into Proposition~\ref{prop:canonical-stress-tensor-translations}
gives the displayed \(T^\mu{}_\nu\).  Since
\[
  T^{\mu\nu}
  =
  \partial^\mu\phi\,\partial^\nu\phi
  -
  \eta^{\mu\nu}
  \left(
    \frac12\partial_\rho\phi\,\partial^\rho\phi+V(\phi)
  \right),
\]
the tensor is symmetric.  Finally, with \(\eta_{00}=-1\),
\[
  \partial_\rho\phi\,\partial^\rho\phi
  =
  -(\partial_0\phi)^2+\sum_{i=1}^{D-1}(\partial_i\phi)^2,
\]
and the displayed expression for \(T^{00}\) follows.
\end{proof}

\section{Metric Variation and Improvement}

\begin{definition}[Lorentzian Hilbert stress tensor]
\label{def:lorentzian-hilbert-stress-tensor}
Suppose the flat-space action is obtained by restricting a generally covariant
background-metric action \(S_g[\phi]\) to \(g_{\mu\nu}=\eta_{\mu\nu}\).  With
the Lorentzian mostly-plus action convention used in this volume, the Hilbert
stress tensor is defined by
\[
  T_{\mathrm H}^{\mu\nu}(x)
  =
  \frac{2}{\sqrt{|g|}}
  \frac{\delta S_g}{\delta g_{\mu\nu}(x)}
  \bigg|_{g=\eta}.
\]
Equivalently,
\[
  \delta_g S_g
  =
  \frac12
  \int\dd^Dx\,\sqrt{|g|}\,
  T_{\mathrm H}^{\mu\nu}\,\delta g_{\mu\nu}.
\]
Using
\(\delta g^{\mu\nu}=-g^{\mu\rho}g^{\nu\sigma}\delta g_{\rho\sigma}\), this is
the same as
\[
  T_{\mathrm H}^{\mu\nu}(x)
  =
  -\frac{2}{\sqrt{|g|}}
  g^{\mu\rho}g^{\nu\sigma}
  \frac{\delta S_g}{\delta g^{\rho\sigma}(x)}
  \bigg|_{g=\eta}.
\]
The Euclidean source-functional chapters state their Euclidean sign convention
separately.  The chain-rule conversion between covariant and contravariant
metric variations is the same; the sign attached to the stress tensor depends
on whether one is varying the Lorentzian action or the Euclidean
\(\exp(-S_E)\) source functional.
\end{definition}

\begin{proposition}[Minimal scalar Hilbert tensor]
\label{prop:minimal-scalar-hilbert-stress}
For the minimally coupled scalar action
\[
  S_g[\phi]
  =
  \int\dd^Dx\,\sqrt{|g|}
  \left(
    -\frac12g^{\rho\sigma}\partial_\rho\phi\,\partial_\sigma\phi
    -V(\phi)
  \right),
\]
Definition~\ref{def:lorentzian-hilbert-stress-tensor} gives
\[
  T_{\mathrm H}^{\mu\nu}
  =
  \partial^\mu\phi\,\partial^\nu\phi
  -
  g^{\mu\nu}
  \left(
    \frac12\partial_\rho\phi\,\partial^\rho\phi+V(\phi)
  \right),
\]
which agrees in flat space with
Proposition~\ref{prop:real-scalar-stress-tensor}.
\end{proposition}

\begin{proof}
Vary with respect to the contravariant metric.  The determinant identity is
\[
  \delta\sqrt{|g|}
  =
  -\frac12\sqrt{|g|}\,g_{\mu\nu}\delta g^{\mu\nu}.
\]
Therefore
\[
  \delta_g S_g
  =
  -\frac12
  \int\dd^Dx\,\sqrt{|g|}
  \left[
    \partial_\mu\phi\,\partial_\nu\phi
    +
    g_{\mu\nu}
    \left(
      -\frac12g^{\rho\sigma}\partial_\rho\phi\,\partial_\sigma\phi
      -V
    \right)
  \right]\delta g^{\mu\nu}.
\]
Comparing with the contravariant form of
Definition~\ref{def:lorentzian-hilbert-stress-tensor} gives
\[
  T_{\mathrm H\,\mu\nu}
  =
  \partial_\mu\phi\,\partial_\nu\phi
  +
  g_{\mu\nu}
  \left(
    -\frac12\partial_\rho\phi\,\partial^\rho\phi
    -V
  \right).
\]
Raising both indices gives the displayed expression.
\end{proof}

\begin{definition}[Stress-tensor improvement]
\label{def:stress-tensor-improvement}
Stress tensors that differ by
\[
  T^{\mu\nu}\longmapsto
  T^{\mu\nu}+\partial_\rho B^{\rho\mu\nu},
  \qquad
  B^{\rho\mu\nu}=-B^{\mu\rho\nu},
\]
have the same local conservation law.  They define the same translation
charges when the resulting surface integrals vanish.  The choice of
representative matters for local operator identities and for couplings to
background geometry, so the representative is part of the framework.
\end{definition}

\section{Lorentz Currents}

\begin{proposition}[Lorentz current from a symmetric stress tensor]
\label{prop:lorentz-current-symmetric-stress}
For a scalar field, an infinitesimal Lorentz transformation is specified by
\(\omega_{\mu\nu}=-\omega_{\nu\mu}\) and
\[
  \delta_\omega\phi(x)
  =
  -\omega^\mu{}_\nu x^\nu\partial_\mu\phi(x).
\]
If \(T^{\mu\nu}=T^{\nu\mu}\) and \(\partial_\mu T^{\mu\nu}=0\), then
\[
  M^{\mu,\rho\sigma}
  =
  x^\rho T^{\mu\sigma}-x^\sigma T^{\mu\rho}
\]
is conserved:
\[
  \partial_\mu M^{\mu,\rho\sigma}=0.
\]
The Lorentz current for parameter \(\omega_{\rho\sigma}\) is
\[
  j^\mu_\omega
  =
  \frac12\omega_{\rho\sigma}M^{\mu,\rho\sigma},
\]
and the corresponding charges, when the spatial integrals exist, are
\[
  J^{\rho\sigma}
  =
  \int\dd^{D-1}\vec x\,
  \left(
    x^\rho T^{0\sigma}-x^\sigma T^{0\rho}
  \right).
\]
\end{proposition}

\begin{proof}
Differentiate \(M^{\mu,\rho\sigma}\):
\[
  \partial_\mu M^{\mu,\rho\sigma}
  =
  T^{\rho\sigma}+x^\rho\partial_\mu T^{\mu\sigma}
  -
  T^{\sigma\rho}-x^\sigma\partial_\mu T^{\mu\rho}.
\]
Stress-tensor conservation removes the terms proportional to \(x\), and
symmetry of \(T^{\mu\nu}\) removes the remaining difference.  The current
\(j^\mu_\omega\) is therefore conserved for every constant antisymmetric
\(\omega_{\rho\sigma}\).
\end{proof}

\begin{remark}[Spin terms and Belinfante improvement]
\label{rem:spin-terms-belinfante}
For fields carrying Lorentz or spinor indices, the Lorentz current includes an
additional spin current.  In such cases one either keeps the canonical spin
current explicitly or passes to a Belinfante-improved symmetric stress tensor,
provided the improvement is defined and the boundary conditions make the
charges equivalent.  The spinor chapter fixes the corresponding representation
matrices before using these currents.
\end{remark}

\section{Quantum Currents and Charges}

\begin{definition}[Quantum current]
\label{def:quantum-current-operator-valued-distribution}
In a quantum field theory, a current is a vector of operator-valued
distributions \(j^\mu_R\) on a common dense domain
\(\mathcal D\subset\Hilb\).  For \(f\in C_c^\infty(\mathbb M^D)\), the
smeared current is denoted
\[
  j^\mu_R(f)=\int\dd^Dx\,f(x)j^\mu_R(x).
\]
Current conservation is a distributional statement:
\[
  \partial_\mu j_R^\mu=0.
\]
Equivalently, for all \(\psi,\chi\in\mathcal D\) and all
\(f\in C_c^\infty(\mathbb M^D)\),
\[
  \sum_{\mu=0}^{D-1}
  \left\langle\psi,j_R^\mu(\partial_\mu f)\chi\right\rangle=0,
\]
with the derivative acting on the test function.  This is the operator-valued
distributional version of integration by parts.
\end{definition}

\begin{definition}[Regulated Ward identity and contact terms]
\label{def:regulated-ward-identity-contact-terms}
In correlation functions with local insertions, the quantum statement is a
Ward identity of distributions and includes contact terms on collision
diagonals.  Its precise form is part of the chosen quantum construction.  In a
regulated source-functional presentation, let
\[
  X=\prod_{i=1}^n\mathcal O_i(x_i)
\]
be a product of renormalized local insertions and let the local symmetry
parameter be \(\epsilon(x)\).  The finite-regulator change of variables gives
\begin{equation}
  \int\dd^Dx\,(\partial_\mu\epsilon)(x)\,
  \left\langle j_R^\mu(x)X\right\rangle_{\rm reg}
  =
  \left\langle \delta_{\epsilon R}X\right\rangle_{\rm reg}
  +
  \left\langle \mathcal A_{\epsilon,R}X\right\rangle_{\rm reg}.
  \label{eq:quantum-noether-ward-regulated}
\end{equation}
Here \(\mathcal A_{\epsilon,R}\) is the local term produced by the variation
of the regulated integration density and by the local counterterms in the
source chart; in a fermionic path integral this density is the finite
Berezinian density described in
Chapter~\ref{ch:spinor-fields-fermionic-path-integrals}.  If the symmetry is
implemented without anomaly in the chosen regularization and source chart,
\(\mathcal A_{\epsilon,R}=0\).  The first term on the right of
\eqref{eq:quantum-noether-ward-regulated} gives contact distributions
supported at \(x=x_i\):
\[
  \delta_{\epsilon R}X
  =
  \sum_i
  \epsilon(x_i)\,
  \mathcal O_1(x_1)\cdots
  (\delta_R\mathcal O_i)(x_i)\cdots
  \mathcal O_n(x_n),
\]
with the corresponding spin or derivative-of-delta terms when the inserted
operators carry spacetime indices or derivatives.  After the regulator is
removed, \eqref{eq:quantum-noether-ward-regulated} becomes an identity in the
renormalized source-functional contact-term chart.  The fixed-point version
used later is formulated in
Section~\ref{sec:cft-source-functionals-ward-contacts}, with the status of
the source brackets specified in
Definition~\ref{def:cft-fixed-point-source-bracket-status}.
\end{definition}

\begin{definition}[Quantum charge as a smeared limit]
\label{def:quantum-charge-smeared-limit}
A quantum charge is defined as a limit of smeared local charges when that
limit exists.  Choose a compactly supported time smearing \(f_t(x^0)\)
concentrated near \(t\), normalized by
\(\int\dd x^0\,f_t(x^0)=1\), and spatial cutoff functions
\(\chi_R(\vec x)\) with \(\chi_R\to1\) as \(R\to\infty\).  The regulated
charge is
\[
  Q_R[f_t,\chi_R]
  =
  \int\dd^Dx\,f_t(x^0)\chi_R(\vec x)j_R^0(x).
\]
If these operators converge on a specified domain, the limit is the charge
operator \(Q_R\).  If \(Q_R\) is essentially self-adjoint or has a chosen
self-adjoint extension, it exponentiates to a one-parameter unitary group by
Stone's theorem.
\end{definition}

\begin{definition}[Quantum generator convention]
\label{def:quantum-generator-convention}
With the fixed-coordinate Noether convention, the infinitesimal quantum action
on an operator \(\widehat{\mathcal O}\) is
\[
  \delta_R\widehat{\mathcal O}
  =
  \frac{\ii}{\hbar}
  [Q_R,\widehat{\mathcal O}]
\]
on the common domain where the commutator is defined.  For spacetime
translations the covariance convention of the earlier chapters is
\[
  U(a)^{-1}\widehat\Phi_A(x)U(a)
  =
  \widehat\Phi_A(x+a),
  \qquad
  U(a)=\exp\left(\frac{\ii}{\hbar}a^\nu P_\nu\right).
\]
\end{definition}

\begin{proposition}[Translation commutator convention]
\label{prop:translation-commutator-convention}
The convention in Definition~\ref{def:quantum-generator-convention} is
equivalent to
\[
  [P_\nu,\widehat\Phi_A(x)]
  =
  \ii\hbar\,\partial_\nu\widehat\Phi_A(x)
\]
in distributional notation.
\end{proposition}

\begin{proof}
Differentiate
\[
  U(a)^{-1}\widehat\Phi_A(x)U(a)=\widehat\Phi_A(x+a)
\]
with respect to \(a^\nu\) at \(a=0\), using
\(U(a)=\exp(\ii a^\rho P_\rho/\hbar)\).  The derivative of the left hand side
is
\[
  -\frac{\ii}{\hbar}P_\nu\widehat\Phi_A(x)
  +
  \frac{\ii}{\hbar}\widehat\Phi_A(x)P_\nu
  =
  -\frac{\ii}{\hbar}[P_\nu,\widehat\Phi_A(x)].
\]
The derivative of the right hand side is
\(\partial_\nu\widehat\Phi_A(x)\).  Multiplying by \(\ii\hbar\) gives the
claimed commutator.
\end{proof}

\begin{definition}[Stress-tensor charges in the quantum theory]
\label{def:quantum-stress-tensor-charges}
When a renormalized stress tensor exists as a local operator-valued
distribution, the Poincare generators are represented by
\[
  P^\nu=\int\dd^{D-1}\vec x\,T^{0\nu}(x),
\]
and, for a symmetric stress tensor in the scalar case,
\[
  J^{\rho\sigma}
  =
  \int\dd^{D-1}\vec x\,
  \left(
    x^\rho T^{0\sigma}(x)-x^\sigma T^{0\rho}(x)
  \right),
\]
with the same regulator-and-limit interpretation as for general charges.
The quantum theory must specify the operator domains, renormalization
prescription, conservation law, and boundary behavior under which the
classical Noether densities define these charges.
\end{definition}

Once those data are fixed, the stress tensor is the local bridge between
spacetime symmetry and the generators that act on states and fields.
